\vspace{-0.3cm}
Many biological processes can be accurately described by partial differential equations (PDEs). In particular, reaction-diffusion equations can be used to model the interactions in a system of chemical substances spreading through a domain and reacting with each other. In this thesis, we study various models for a biological process related to cell communication.

Cells can communicate with one another by producing signalling molecules referred to as \emph{ligands} and releasing them into their environment. These ligands travel through the extracellular space until they are captured by a \emph{receptor} on the membrane of another cell, forming a bounded receptor, or ligand-receptor \emph{complex}. Specifically, we consider the case where these receptors couple to a \emph{G-protein}. These receptors can diffuse across the cell membrane. When a ligand binds to a receptor, a reaction is triggered that passes the signal through the membrane and activates a G-protein, which further passes the signal to the interior of the cell, where it can be processed appropriately. Eventually, the ligand is released by the receptor, enabling the ligand and receptor to repeat the process of diffusing across their domains and reacting with other receptors and ligands, respectively.

Due to their diffusing nature, one might expect to find G-protein coupled receptors homogeneously distributed across the cell membrane. However, experiments have shown that this is not the case. Instead, they form clusters of various sizes, which can influence their signalling behaviour \cite{kockerkoren24,mandala25}. We aim to formulate and study reaction-diffusion systems predicting such clustering using bulk-surface models, since receptors are confined to the cell membranes.

In this thesis, we study various systems of bulk-surface systems modelling the concentrations of ligands, receptors and their complex. \cref{ch:Preliminaries} establishes the required mathematical background in the field of PDEs, introducing notions such as weak derivatives and weak solutions to PDEs, as well as recalling some simple but important inequalities that are essential in later proofs. \cref{ch:Bulk_surface_model} derives a model involving ordinary differential equations (ODEs) for the ligand-receptor interactions with assumptions from \cite{iber2014} and discusses a classical proof of existence and uniqueness of weak solutions to his model, as well as some qualitative properties. \cref{ch:Turing_patterns} introduces the theory of Turing patterns, which is often used to account for spatial organisation in mathematical biology. \cref{ch:Numerical_Analysis} introduces the finite element method to simulate another variant of model. Using these simulations, we seek to explain the formation of receptor clusters by this model's ability to generate Turing patterns. We also reduce the model to a system of fewer components and analyse how well this reduced variant represents the full system. 

This thesis is written under the assumption that the reader is comfortable with mathematical theory typically taught in a mathematics bachelor's programme. In particular, relative fluency in subjects such as multivariate calculus, linear algebra and ODEs is required. Knowledge on the fundamental definitions and theory of PDEs, in particular the diffusion equation, and functional analysis is useful.